\section{Sparse Specification Search}

% Scope:
% Collect the discrete-support search results in one place. This file should
% cover the sparse menu of conventional regression adjustments, with subheadings
% separating the regimes and the studentized consequences.

\subsection{Roadmap}

% Explain the discrete search geometry: the analyst chooses a subset of controls
% from a combinatorial menu, so the extreme value depends on both the number of
% admissible supports and their dependence structure.
%
% Clarify why the regimes are separated:
% low-dimensional search s L_{s,p} = o(n) allows a sharp linearized/root-process
% approximation; critical search n \asymp s L_{s,p} requires upper/lower
% comparison bounds; studentized search is what connects the coefficient
% geometry to significance and replication claims.
%
% This subsection should be mostly explanatory, with a small summary table or
% paragraph listing the main results proved below.

\subsection{Below the Critical Regime}

% Main purpose:
% show that when s L_{s,p} = o(n), the original p-hacking/root problem is close
% to a simpler order-zero or linearized surrogate process.
%
% Include:
% - the root formulation and surrogate root;
% - assumptions for sparse overlap, covariance-normalized Gram concentration,
%   and drift control;
% - the main comparison theorem and its rate;
% - interpretation of the assumptions in transparent special cases such as
%   d = 0, beta_{-0} = 0, and diagonal C.

\subsection{The Critical Sparse Regime}

% Main purpose:
% characterize the order of sup_{S in S_s} \hat beta_{0,S} when
% n \asymp s L_{s,p}.
%
% Include:
% - upper comparison around the deterministic sparse center Theta_{C,s};
% - lower comparison under richness/lower-tail/extraction assumptions;
% - simplified corollaries for no confounding, no control signal, and the
%   baseline null;
% - a clear explanation of what is primitive, what is conditional on technical
%   comparison lemmas, and where the lower-bound assumptions are verified.

\subsection{Studentized Search and Replicability}

% Main purpose:
% translate coefficient-scale search into the statistic actually used for
% significance claims, and explain how studentization changes the geometry.
%
% Include:
% - critical-regime upper and lower transfers for sup \hat t_{0,S};
% - low-dimensional numerator linearization and Gaussian approximation;
% - denominator-flatness conditions and their population interpretations;
% - practical consequences for false positives, threshold crossing, and
%   replication/robustness claims.
%
% Make clear when a t-statistic result is merely a transfer from coefficient
% bounds and when studentization creates a genuinely different optimization
% problem.
